\documentclass[tikz,border=10pt]{standalone}
\usepackage{ctex}
\usepackage{amsmath}
\usetikzlibrary{calc, arrows.meta, shapes.geometric, positioning, backgrounds}

\begin{document}
% 向量综合示意图：向量条件 + 几何曲线 + 代数运算的交叉关系
\begin{tikzpicture}[
  font=\small,
  oval/.style={draw=#1!70, fill=#1!12, ellipse,
               minimum width=3.2cm, minimum height=1.4cm,
               align=center, font=\small},
  arr/.style={-Stealth, thick, #1},
  note/.style={font=\scriptsize, align=center, text width=2.8cm},
]

%─── 三个椭圆（交叉核心） ──────────────────────────
\node[oval=blue]    (vec)  at (3.0, 5.5) {\textbf{向量条件}\\$\vec{OA}\cdot\vec{OB}=0$\\$|\vec{PA}|=|\vec{PB}|$};

\node[oval=orange]  (geom) at (0.2, 1.8) {\textbf{几何曲线}\\椭圆/双曲线\\抛物线/圆};

\node[oval=green!60!black] (alg) at (5.8, 1.8)
  {\textbf{代数运算}\\三角/函数\\不等式/最值};

%─── 三条双向箭头 ───────────────────────────────────
% 向量 → 几何
\draw[arr=blue!70, bend right=18]
  (vec.south west) to
  node[left, note] {向量表示\\几何关系\\（垂直/等距）}
  (geom.north);

% 向量 → 代数
\draw[arr=orange!80, bend left=18]
  (vec.south east) to
  node[right, note] {数量积\\$=$三角余弦\\$=$ 模长关系}
  (alg.north);

% 几何 ↔ 代数
\draw[arr=green!60!black, bend right=10]
  (geom.east) to
  node[below, note, text width=3.2cm] {曲线方程联立\\代数式求最值}
  (alg.west);
\draw[arr=green!60!black!50, bend right=10, dashed]
  (alg.west) to (geom.east);

%─── 中心标注 ───────────────────────────────────────
\node[draw=gray!60, fill=gray!10, rounded corners=5pt,
      font=\scriptsize\bfseries, align=center, inner sep=5pt]
  at (3.0, 3.3) {向量综合题核心\\三域联动};

%─── 典型交叉场景（底部） ───────────────────────────
\node[draw=blue!40, fill=blue!5, rounded corners=4pt,
      font=\scriptsize, align=left, inner sep=5pt,
      text width=3.0cm] at (0.2, 0.3)
  {\textbf{向量＋曲线}\\
   $\vec{OA}\!\cdot\!\vec{OB}=0$\\$A,B$ 在椭圆上\\
   → 求点轨迹};

\node[draw=orange!50, fill=orange!5, rounded corners=4pt,
      font=\scriptsize, align=left, inner sep=5pt,
      text width=3.0cm] at (5.8, 0.3)
  {\textbf{向量＋三角}\\
   $\vec{a}\!\cdot\!\vec{b}=|\vec{a}||\vec{b}|\cos\theta$\\
   → 恒等式证明\\
   → 最值问题};

%─── 标题 ───────────────────────────────────────────
\node[font=\normalsize\bfseries] at (3.0, 7.0) {向量综合：三域交叉结构图};

\end{tikzpicture}
\end{document}
